\PassOptionsToPackage{table,xcdraw}{xcolor}
\documentclass[11pt, a4paper]{article}

\usepackage[T1]{fontenc}
\usepackage[utf8]{inputenc}
\usepackage{tgpagella}
\usepackage[spanish, es-noshorthands]{babel}
\usepackage{amsmath, amssymb}
\usepackage{booktabs}
\usepackage{tabularx}
\usepackage[most]{tcolorbox}
\usepackage[margin=2.5cm]{geometry}
\usepackage{fancyhdr}

\pagestyle{fancy}
\fancyhf{}
\setlength{\headheight}{16pt}
\lhead{\textbf{UCB Sede Tarija} | IMT-221}
\chead{Hoja Conceptual: Op-Amp Real}
\rhead{Semana 8 - Sesión 2}
\renewcommand{\headrulewidth}{0.4pt}
\definecolor{ucbblue}{RGB}{0, 51, 102}

\begin{document}

\begin{center}
    \Large\textbf{\textcolor{ucbblue}{Conceptos Compactos: Op-Amp Ideal vs Real}}[cite: 12]
\end{center}

\begin{tcolorbox}[colback=yellow!5, colframe=ucbblue, title=\textbf{Regla Diagnóstica de Oro}]
    \textbf{No culpe al Op-Amp} antes de revisar exhaustivamente: alimentación, tierra, pinout, lazo de realimentación, resistores, señales de entrada y la conexión de sus instrumentos[cite: 12].
\end{tcolorbox}

\section*{Tabla de Correspondencias[cite: 12]}
\begin{table}[h]
    \centering
    \renewcommand{\arraystretch}{1.5}
    % Se añade \raggedright\arraybackslash para evitar el Underfull \hbox
    \begin{tabularx}{\textwidth}{>{\raggedright\arraybackslash\hsize=0.8\hsize}X >{\raggedright\arraybackslash\hsize=0.8\hsize}X >{\raggedright\arraybackslash\hsize=1.4\hsize}X}
    \toprule
    \rowcolor[HTML]{EFEFEF}
    \textbf{Hipótesis Ideal} & \textbf{Desviación Real} & \textbf{Consecuencia Observable} \\ \midrule
    $V_+ \approx V_-$ bajo lazo cerrado & $V_{OS}$ (Input Offset Volt.) & Offset de salida residual DC[cite: 12] \\
    $i_+ = i_- = 0$ & $I_B$ (Input Bias Curr.) & Error dependiente de la resistencia[cite: 12] \\
    Corrientes idénticas & $I_{OS}$ (Input Offset Curr.) & Error de desbalance residual[cite: 12] \\
    Ratios de resistores exactos & Tolerancia / Mismatch & Resta imperfecta / Error de ganancia[cite: 12] \\
    Predicción de salida exacta & Efectos combinados & Discrepancia Teoría-Medición[cite: 12] \\
    \bottomrule
    \end{tabularx}
\end{table}

\section*{Modelos Analíticos de No Idealidades DC[cite: 12]}
\textbf{1. Tensión de Offset de Entrada ($V_{OS}$)[cite: 12]} \\
Pequeño error diferencial equivalente referido a la entrada. Para el amplificador diferencial acoplado de Lab 3, el error aproximado propagado a la salida es:
\begin{equation*}
    V_{o,OS} \approx \left(1 + \frac{R_2}{R_1}\right)V_{OS} = A_N V_{OS} \text{[cite: 12]}
\end{equation*}

\textbf{2. Corrientes de Polarización y de Offset ($I_B$ e $I_{OS}$)[cite: 12]} \\
Las corrientes generan caídas de tensión parásitas al atravesar las resistencias equivalentes vistas por los terminales:
\begin{equation*}
    V_{err} \approx I_B R_{eq} \text{[cite: 12]}
\end{equation*}
La corriente de offset es la diferencia neta entre ambas entradas:
\begin{equation*}
    I_{OS} = |I_{B+} - I_{B-}| \text{[cite: 12]}
\end{equation*}

\textbf{3. Desajuste de Resistores (Mismatch)[cite: 12]} \\
La resta perfecta asume matemáticamente que $\frac{R_2}{R_1} = \frac{R_4}{R_3}$. Las desviaciones introducen coeficientes desiguales que filtran voltaje en modo común a la salida[cite: 12].

\end{document}
